\documentclass[11pt,a4paper]{article}
\usepackage[margin=25mm,headheight=15pt]{geometry}
\usepackage[T1]{fontenc}
\usepackage[utf8]{inputenc}
\usepackage{lmodern,microtype}
\usepackage{amsmath,amssymb,amsthm,mathtools}
\usepackage{booktabs,longtable,array}
\usepackage{xcolor,listings,fancyhdr,xurl}
\usepackage[colorlinks=true,linkcolor=blue!45!black,citecolor=green!30!black,urlcolor=blue!45!black]{hyperref}
\hypersetup{pdftitle={Real-branch asymptotic inversion at a singular endpoint},pdfauthor={Research and implementation report prepared for Vladimir Reshetnikov},pdfsubject={Power-logarithmic reversion with irrational exponents and a Wolfram Language package}}
\pagestyle{fancy}\fancyhf{}
\fancyhead[L]{\small Real-branch asymptotic inversion}
\fancyhead[R]{\small Power--logarithmic reversion}
\fancyfoot[C]{\thepage}
\setlength{\parskip}{0.35em}
\setlength{\parindent}{1.1em}
\setlength{\emergencystretch}{3em}
\allowdisplaybreaks
\numberwithin{equation}{section}
\newtheorem{theorem}{Theorem}[section]
\newtheorem{proposition}[theorem]{Proposition}
\newtheorem{lemma}[theorem]{Lemma}
\newtheorem{corollary}[theorem]{Corollary}
\theoremstyle{definition}
\newtheorem{definition}[theorem]{Definition}
\newtheorem{example}[theorem]{Example}
\theoremstyle{remark}
\newtheorem{remark}[theorem]{Remark}
\newcommand{\R}{\mathbb R}
\newcommand{\N}{\mathbb N_0}
\newcommand{\OO}{\mathcal O}
\newcommand{\dd}{\mathrm d}
\newcommand{\D}{\mathcal D}
\newcommand{\val}{\operatorname{val}}
\newcommand{\Res}{\operatorname{Res}}
\newcommand{\supp}{\operatorname{supp}}
\newcommand{\abs}[1]{\left|#1\right|}
\newcommand{\norm}[1]{\left\|#1\right\|}
\lstset{basicstyle=\ttfamily\small,columns=fullflexible,keepspaces=true,
 breaklines=true,showstringspaces=false,frame=single,
 rulecolor=\color{black!25},backgroundcolor=\color{black!2},
 aboveskip=0.8em,belowskip=0.8em}

\begin{document}
\begin{titlepage}
\centering
\vspace*{16mm}
{\LARGE\bfseries Real-branch asymptotic inversion\\[0.3em]at a singular endpoint\par}
\vspace{8mm}
{\Large Power--logarithmic expansions with irrational exponents,\\
convergent reversion, and a Wolfram Language implementation\par}
\vspace{12mm}
{\large Research and implementation report\\
Prepared for Vladimir Reshetnikov\\[0.4em]7 September 2026\par}
\vspace{12mm}
\begin{minipage}{0.91\textwidth}
\begin{abstract}
The inverse of $x+x^2(1+\log x)$ on the positive real axis has logarithmic
coefficient blocks, not a Taylor series at zero. The inverse of
$x+x^{\sqrt2}$ has irrational exponents and does not fit a rational Puiseux
grid. We solve both problems within a common constructive theory for
\[
 f(x)=a x^p\left(1+\sum_{i=1}^{m}x^{\delta_i}B_i(\log x)\right),
 \qquad a,p,\delta_i>0,\quad B_i\in\R[L].
\]
An explicit finite differential-operator formula computes every inverse
block. The inverse support lies in a positive finitely generated additive
semigroup, so every power cutoff requires finitely many multi-indices even
when the exponent group is dense. We prove positive-branch existence and
uniqueness, the coefficient formula, actual convergence for each finite
model sufficiently near zero, a finite-boundary remainder theorem, and
transport of differentiated forward remainders. Resonances, logarithmic
orders, and nonunit leading powers are treated explicitly.

The accompanying package implements exact algebraic-exponent enumeration,
polynomial logarithmic coefficients, complete-block truncation, guarded
input precision, and an independent sparse-composition residual check.
A separate marker-based reversion function handles general near-identity
perturbations without pretending that marker order is automatically
asymptotic order. All limitations and verification status are explicit:
the mathematical formulas have independent exact and high-precision
checks, but native Wolfram-kernel execution was unavailable in this session.
\end{abstract}
\end{minipage}
\vfill
{\small Package version 1.0.0 \quad\textbullet\quad Source, examples, tests, and verification scripts included\par}
\end{titlepage}

\tableofcontents
\clearpage

\section{The problem and the answer}
\label{sec:answer}

The question supplied with this task, and its public version
\cite{question}, asks for the positive real inverse as the argument tends to
zero. Its two examples are
\begin{equation}
 f_1(x)=x+x^2(1+\log x),\qquad
 f_2(x)=x+x^{\sqrt2},\qquad x>0.
\end{equation}
The archived question reports unsuccessful attempts involving
\texttt{InverseFunction}, \texttt{Series}, \texttt{Asymptotic}, and
\texttt{InverseSeries}. Those are historical observations, not tests of a
current Wolfram release performed for this report.

Put $L=\log y$. The first inverse is
\begin{align}
 g_1(y)={}&y-y^2(1+L)+y^3(2L^2+5L+3)\nonumber\\
 &-y^4\left(5L^3+\frac{41}{2}L^2+27L+\frac{23}{2}\right)\nonumber\\
 &+y^5\left(14L^4+\frac{241}{3}L^3+\frac{335}{2}L^2
                      +151L+\frac{299}{6}\right)
 +\OO\left(y^6(1+|L|)^5\right).
 \label{eq:first-answer}
\end{align}
The second inverse is
\begin{align}
 g_2(y)={}&y-y^{\sqrt2}+\sqrt2\,y^{2\sqrt2-1}
 -\frac{6-\sqrt2}{2}\,y^{3\sqrt2-2}
 +\OO\left(y^{4\sqrt2-3}\right).
 \label{eq:second-answer}
\end{align}
Every limit and remainder in these displays is for $y\to0^+$.
Neither formula requires selecting a complex branch of an unspecified
\texttt{InverseFunction} object.

\paragraph{Why logarithms must remain in the remainder.}
After retaining only $y-y^2(1+\log y)$, the error is
\begin{equation}
 y^3(2\log^2y+5\log y+3)
 +\OO\bigl(y^4(1+|\log y|)^3\bigr).
 \label{eq:not-o-y3}
\end{equation}
It is therefore \emph{not} $\OO(y^3)$ in the usual analytic meaning of
big-O. A valid bound is $\OO(y^3(1+|\log y|)^2)$. Logarithms may occur inside
Wolfram's series coefficients \cite{wlseriesdata,wlseries}; this does not
license dropping them from an analytic remainder bound.

\paragraph{Minimal package use.}
After extracting the archive, load the file with an appropriate local path:
\begin{lstlisting}
Get["RealInverseAsymptotics/Kernel/RealInverseAsymptotics.wl"];
Clear[x, y];
r = RealInverseAsymptotic[
  x + x^2 (1 + Log[x]), {x, 0}, {y, 6}];
Normal[r]
r["Remainder"]
InverseResidual[r]["ZeroBelowCutoff"]
\end{lstlisting}
The cutoff is \emph{exclusive}: \texttt{\{y,6\}} retains every complete
logarithmic block with power smaller than $6$. The predicted outputs are
\eqref{eq:first-answer} without its remainder,
\texttt{PowerLogRemainder[y,6,5]}, and \texttt{True}. These are mathematical
expectations encoded in the regression suite, not a claim of a native-kernel
run in this session. Section~\ref{sec:verification} records what was run.

\section{Selecting a real branch before expanding}
\label{sec:branch}

\subsection{The first example is globally invertible on the positive axis}

\begin{proposition}
The map $f_1:(0,\infty)\to(0,\infty)$ is a strictly increasing bijection.
Its inverse satisfies $g_1(y)/y\to1$ as $y\to0^+$.
\end{proposition}
\begin{proof}
Differentiation gives
\[
 f_1'(x)=1+3x+2x\log x.
\]
The function $3x+2x\log x$ has derivative $5+2\log x$, so its global
minimum occurs at $x=\exp(-5/2)$ and is $-2\exp(-5/2)$. Consequently
\begin{equation}
 f_1'(x)\ge m_0:=1-2\exp(-5/2)>0\qquad(x>0).
 \label{eq:global-slope}
\end{equation}
Also $f_1(x)\to0$ at zero and $f_1(x)\to\infty$ at infinity. Strict
monotonicity and the intermediate value theorem prove the bijection.
Finally $f_1(x)/x=1+x(1+\log x)\to1$, and substitution $x=g_1(y)$ proves
the inverse equivalence.
\end{proof}

The continuous extension has $f_1(0)=0$ and a right derivative equal to one,
but
\[
 f_1''(x)=5+2\log x\longrightarrow-\infty.
\]
Likewise $g_1''(y)=-f_1''(g_1(y))/f_1'(g_1(y))^3\to+\infty$.
Thus there is no ordinary second-order Taylor expansion of the smooth-at-zero
kind that an analytic inverse theorem would require. A nonzero limiting
first derivative alone is not analyticity.

The constant in the historical complex output has the form
\[
 c=\exp\bigl(-1+W_{-1}(-e)\bigr).
\]
The nonzero zeros of the complex continuation formally satisfy
$1+x(1+\log x)=0$, which becomes $t e^t=-e$ after $t=1+\log x$.
A local expansion about such a nonzero complex preimage is a different
problem from the real branch characterized by $g_1(y)\to0^+$. Reality of
some symbols is not a substitute for specifying the limiting preimage and
the domain of the inverse. The package never asks \texttt{InverseFunction}
to choose that branch; it constructs the positive local branch directly.
The documented role of \texttt{InverseFunction} and the series-reversion
role of \texttt{InverseSeries} are distinct \cite{wlinversefunction,wlinverseseries}.

\subsection{The general branch condition}

Our central model is
\begin{equation}
 f(x)=a x^p\left(1+\sum_{i=1}^{m}x^{\delta_i}B_i(\log x)\right),
 \qquad a>0,\quad p>0,\quad\delta_i>0,
 \label{eq:model}
\end{equation}
where each $B_i$ is a real polynomial and all parameters are fixed.
The empty sum, $m=0$, is permitted. Equal $\delta_i$ should first be
combined; zero blocks may be deleted.

\begin{lemma}[Positive local branch]
For \eqref{eq:model}, there is an $\varepsilon>0$ such that $f$ is positive
and strictly increasing on $(0,\varepsilon)$. Its inverse $g$ on the
corresponding image satisfies
\begin{equation}
 g(y)\sim z,\qquad z=(y/a)^{1/p}>0.
 \label{eq:leading-root}
\end{equation}
This characterizes the branch used throughout the report.
\end{lemma}
\begin{proof}
For every $\delta>0$ and fixed integer $d\ge0$,
$x^\delta(1+|\log x|)^d\to0$. This follows by putting $x=e^{-t}$ and
comparing $e^{-\delta t}(1+t)^d$ with zero as $t\to\infty$.
Differentiating each correction in \eqref{eq:model} shows
\[
 f(x)\sim a x^p,\qquad f'(x)\sim apx^{p-1}>0.
\]
Positivity and monotonicity follow on a sufficiently small interval.
The inverse exists there and tends to zero. Substituting it into the first
equivalence gives $y\sim a g(y)^p$, hence \eqref{eq:leading-root}.
\end{proof}

This lemma is local: a general finite model can turn around away from the
endpoint. Global invertibility is established separately for the two
original examples, not presumed for every package input. Real powers are
always taken on positive arguments, using $x^r=\exp(r\log x)$.
This restriction is essential; Wolfram's general complex \texttt{Power}
uses principal values \cite{wlpower}. The implementation does not apply
\texttt{PowerExpand} to unknown-sign expressions.

\section{The algebra of power--logarithmic blocks}
\label{sec:algebra}

\subsection{Blocks and their actual asymptotic ordering}

Let $L=\log z$ and $\D=\dd/\dd L$. A block is $z^\alpha P(L)$, with
$P$ a polynomial. If $\alpha<\beta$ and $P,Q$ are nonzero polynomials,
\begin{equation}
 z^\beta Q(\log z)=o\bigl(z^\alpha P(\log z)\bigr).
 \label{eq:dominance}
\end{equation}
Indeed, the ratio has a positive power of $z$ multiplied by a rational
function of $\log z$, and the exponential comparison just used applies.
Within one block, larger logarithmic degree means larger magnitude.
A fully ordered monomial scale therefore orders the power increasingly and,
at equal power, the logarithmic degree decreasingly.

The software retains \emph{whole polynomial blocks}. This avoids treating
$L$ as constant during differentiation or silently discarding a lower
logarithmic term at a retained power. It also gives an unambiguous meaning
to a cutoff. A request for ``four terms'' would otherwise be ambiguous:
there may be four monomials inside just two power blocks.

\begin{lemma}[Differential block identity]
For every nonnegative integer $r$,
\begin{equation}
 \frac{\dd^r}{\dd z^r}\bigl[z^\alpha P(\log z)\bigr]
 =z^{\alpha-r}\prod_{j=0}^{r-1}(\D+\alpha-j)P(L).
 \label{eq:block-derivative}
\end{equation}
The empty product, for $r=0$, is the identity operator.
\end{lemma}
\begin{proof}
One differentiation gives
$z^{\alpha-1}(\alpha P+P')$. Repeated application proves the result by
induction. The factors commute because each is a constant plus $\D$.
\end{proof}

This identity is the computational bridge from a derivative formula to
finite polynomial arithmetic. A derivative of $P(\log z)$ is never omitted.
For instance,
\[
 \frac12\frac{\dd}{\dd y}\{y^4(1+\log y)^2\}
 =y^3\{2(1+\log y)^2+(1+\log y)\},
\]
which is exactly the third block of \eqref{eq:first-answer}.

\subsection{Irrational exponents do not destroy local finiteness}

For the positive gaps in \eqref{eq:model}, set
\begin{equation}
 S=\left\{n_1\delta_1+\cdots+n_m\delta_m:n_i\in\N\right\}.
 \label{eq:semigroup}
\end{equation}
If $\sum n_i\delta_i<H$, then $n_i<H/\delta_i$ for every $i$.
Consequently there are only finitely many admissible multi-indices below
any finite $H$, and only finitely many support values there.
The additive \emph{group} generated by two incommensurable gaps may be
dense, but it allows negative integer coefficients. The positive semigroup
\eqref{eq:semigroup} does not have that defect.

We can therefore work with formal sums
\[
 U(z)=\sum_{\sigma\in S\setminus\{0\}}z^\sigma C_\sigma(L),
 \qquad C_\sigma\in\R[L],
\]
in a completion with respect to increasing power order. Products, and the
formal series $\log(1+U)$ and $(1+U)^r$, are defined coefficientwise:
at every fixed cutoff only finitely many products contribute.
The valuation $\val U$ is the least nonzero power in its support;
logarithmic degree is tracked separately. This restricted, constructive
setting is enough here. General transseries and their compositional
inversion live in a substantially larger framework \cite{edgarbegin,edgarcomp}.

\paragraph{Resonances.}
Different multi-indices can have the same weight. If $\delta_2=2\delta_1$,
then $(2,0)$ and $(0,1)$ coincide in weight. Their polynomial coefficients
must be added, not counted as different asymptotic terms. Cancellation can
remove an entire block. Exact algebraic canonicalization in the package is
used precisely for this reason; floating-point ordering or tolerance-based
equality would be an unsound replacement.

\section{Reversion with an auxiliary marker}
\label{sec:marker}

\subsection{A derivative formula that never evaluates the singular endpoint}

The underlying identity is a standard form of Lagrange--B\"urmann inversion
\cite{gessel}. We give its proof to make clear what is analytic, what is
formal, and why logarithms at zero are not an obstruction.

\begin{theorem}[Marker reversion]
Fix $q>0$. Suppose $h$ and $\Phi$ are analytic in a complex neighborhood of
$q$. Let $u(q,\eta)$ be the solution of
\[
 u+\eta h(u)=q,\qquad u(q,0)=q.
\]
Then, for sufficiently small $|\eta|$,
\begin{equation}
 \Phi(u(q,\eta))=\Phi(q)+
 \sum_{N\ge1}\frac{(-\eta)^N}{N!}
 \frac{\dd^{N-1}}{\dd q^{N-1}}
 \left[\Phi'(q)h(q)^N\right].
 \label{eq:marker-theorem}
\end{equation}
It is also an identity of formal power series in the marker.
\end{theorem}
\begin{proof}
Put $u=q+s$, $\phi(s)=-h(q+s)$, and $\Psi(s)=\Phi(q+s)$. Then
$s=\eta\phi(s)$. Initially assume $\phi(0)\ne0$. Changing variable
$\eta=s/\phi(s)$ in the residue for the coefficient of $\eta^N$ gives
\begin{align*}
 [\eta^N]\Psi(s(\eta))
 &=\Res_{s=0}\Psi(s)\left\{
       \frac{\phi(s)^N}{s^{N+1}}
       -\frac{\phi'(s)\phi(s)^{N-1}}{s^N}\right\}\dd s\\
 &=\frac1N\Res_{s=0}\frac{\Psi'(s)\phi(s)^N}{s^N}\dd s
 =\frac1N[s^{N-1}]\Psi'(s)\phi(s)^N.
\end{align*}
The middle equality follows because the residue of
$\dd(\Psi(s)\phi(s)^N/s^N)/\dd s$ is zero. Replacing the coefficient by a
Taylor derivative gives \eqref{eq:marker-theorem}. Each coefficient is a
polynomial in finitely many derivatives of $h$ and $\Phi$, so the identity
extends to $\phi(0)=0$ as well. Analytic existence and convergence near
$\eta=0$ follow independently from the analytic implicit function theorem,
since the derivative with respect to $u$ at $(q,0)$ is one.
\end{proof}

For $\Phi(q)=q$, this becomes the particularly simple rule
\begin{equation}
 g(y)=y+\sum_{N\ge1}\frac{(-1)^N}{N!}
 \frac{\dd^{N-1}}{\dd y^{N-1}}h(y)^N,
 \qquad f(x)=x+h(x).
 \label{eq:near-id}
\end{equation}
In \eqref{eq:near-id}, setting the marker to one needs justification.
Formal positive grading supplies a coefficientwise justification; an
analytic smallness argument is still required to identify the sum with an
actual function. Section~\ref{sec:convergence} supplies both for our finite
models. The distinction is not cosmetic: a marker Taylor series can exist
without being a useful asymptotic expansion at the endpoint of interest.

At each fixed positive $q$, functions such as $q^\alpha$ and $\log q$ are
analytic in a small complex disk avoiding zero and the negative axis.
No derivative at $q=0$ occurs in the theorem. This is the essential reason
that the derivative form works for the question while an ordinary Taylor
inverse at zero is the wrong object.

\subsection{A general exactly invertible dominant core}

A useful generalization explains the connection with the supplied
transseries materials. Suppose
\[
 F(x)=H(x)+R(x),\qquad C=H^{-1}
\]
on a specified local branch. Put $t=H(x)$ and $h(t)=R(C(t))$. Then
$t+\eta h(t)=y$ and \eqref{eq:marker-theorem} gives
\begin{equation}
 F_\eta^{-1}(y)=C(y)+\sum_{N\ge1}\frac{(-\eta)^N}{N!}
 \frac{\dd^{N-1}}{\dd y^{N-1}}
 \left[C'(y)R(C(y))^N\right].
 \label{eq:general-core}
\end{equation}
This is a change of coordinate, not a new Lagrange inversion theorem.
The companion volume supplied in the repository uses precisely this kind
of phase-coordinate and scaled reversion, and explicitly distinguishes
formal marker identities from convergence at the physical marker
\cite{repocompanion}. Here the core is $H(x)=a x^p$, and the positive power
gaps provide a finite, directly implementable grading.

\section{The all-order coefficient formula}
\label{sec:master}

For a multi-index $\boldsymbol n=(n_1,\ldots,n_m)\in\N^m$, write
\begin{equation}
 N=|\boldsymbol n|=\sum_i n_i,\qquad
 \boldsymbol n!=\prod_i n_i!,\qquad
 \sigma(\boldsymbol n)=\sum_i n_i\delta_i,\qquad
 B^{\boldsymbol n}(L)=\prod_i B_i(L)^{n_i}.
 \label{eq:multiindex}
\end{equation}
For $N\ge1$, define
\begin{equation}
 C_{\boldsymbol n}(L)=
 \frac{(-1)^N}{p^N\boldsymbol n!}
 \prod_{j=1}^{N-1}\bigl(\D+1+\sigma(\boldsymbol n)+pj\bigr)
 B^{\boldsymbol n}(L).
 \label{eq:coefficient}
\end{equation}
Again the operator product is empty when $N=1$.

\begin{theorem}[Constructive inverse formula]
For the positive branch of \eqref{eq:model}, with $z=(y/a)^{1/p}$ and
$L=\log z$, the inverse has the expansion
\begin{equation}
 \boxed{\displaystyle
 g(y)=z+\sum_{\boldsymbol n\ne0}
          z^{1+\sigma(\boldsymbol n)}C_{\boldsymbol n}(L).}
 \label{eq:master}
\end{equation}
Equal weights are combined by summing their $C_{\boldsymbol n}$.
The formula is a well-defined formal identity. For each fixed finite model,
it also converges absolutely for sufficiently small positive $y$ in the
sense established in Section~\ref{sec:convergence}. In particular, every
finite cutoff is a genuine asymptotic expansion of the specified real
inverse.
\end{theorem}
\begin{proof}[Proof of the coefficient identity]
Set $t=x^p$ and $q=y/a$. Equation \eqref{eq:model} becomes
\[
 q=t+\sum_i t^{1+\delta_i/p} B_i\left(\frac{\log t}{p}\right).
\]
Insert a separate marker for each summand, apply
\eqref{eq:marker-theorem} to $\Phi(t)=t^{1/p}$, and extract the indicated
multinomial coefficient. The coefficient belonging to $\boldsymbol n$ is
\begin{equation}
 \frac{(-1)^N}{\boldsymbol n!}
 \frac{\dd^{N-1}}{\dd q^{N-1}}
 \left[\frac1p q^{N-1+(1+\sigma)/p}
 B^{\boldsymbol n}\left(\frac{\log q}{p}\right)\right].
 \label{eq:raw-coefficient}
\end{equation}
Applying \eqref{eq:block-derivative}, with
$\dd/\dd(\log q)=\D/p$, leaves the power
$q^{(1+\sigma)/p}=z^{1+\sigma}$. Each of the $N-1$ derivative factors
contributes a factor $1/p$, in addition to the displayed $1/p$.
The resulting factors are exactly
$\D+1+\sigma+p,\ldots,\D+1+\sigma+p(N-1)$, proving
\eqref{eq:coefficient}. Local finiteness of \eqref{eq:semigroup} makes every
coefficient and every cutoff finite, so the marker substitution is formally
legitimate. Analytic convergence is proved below without presuming it from
this formal calculation.
\end{proof}

\begin{corollary}[Degree and support bounds]
If $d_i=\deg B_i$, then
\begin{equation}
 \deg C_{\boldsymbol n}\le D(\boldsymbol n):=\sum_i n_i d_i.
 \label{eq:degree-bound}
\end{equation}
Every inverse power belongs to
$\{(1+\sigma)/p:\sigma\in S\}$. Cancellations can remove powers but cannot
create a power outside this set.
\end{corollary}
\begin{proof}
The initial product has degree $D(\boldsymbol n)$; differentiation and
addition of a scalar multiple do not increase polynomial degree.
The support statement is immediate from \eqref{eq:master}.
\end{proof}

\paragraph{First and second corrections.}
For a single block $x^\delta B(L)$, the first two corrections are
\begin{equation}
 g(y)=z-\frac{z^{1+\delta}}pB(L)
 +\frac{z^{1+2\delta}}{2p^2}
 \left\{(1+2\delta+p)B(L)^2+2B(L)B'(L)\right\}+\cdots.
 \label{eq:first-two-general}
\end{equation}
For two different blocks, the multi-index $e_i+e_j$ gives
\begin{equation}
 \frac{z^{1+\delta_i+\delta_j}}{p^2}
 \bigl(\D+1+\delta_i+\delta_j+p\bigr)(B_iB_j).
 \label{eq:cross-term}
\end{equation}
These cross terms are why simply adding separately computed inverse
corrections is incorrect.

\subsection{Powers and the logarithm of the inverse}

The same proof with $\Phi(t)=t^{r/p}$, for any fixed real $r$, yields
\begin{equation}
 g(y)^r=z^r+\sum_{\boldsymbol n\ne0}
 \frac{r(-1)^N z^{r+\sigma}}{p^N\boldsymbol n!}
 \prod_{j=1}^{N-1}(\D+r+\sigma+pj)B^{\boldsymbol n}(L).
 \label{eq:observable-power}
\end{equation}
For $r=0$ the expression is identically one. Differentiating with respect to
$r$ at zero gives
\begin{equation}
 \log g(y)=\log z+\sum_{\boldsymbol n\ne0}
 \frac{(-1)^N z^\sigma}{p^N\boldsymbol n!}
 \prod_{j=1}^{N-1}(\D+\sigma+pj)B^{\boldsymbol n}(L).
 \label{eq:log-inverse}
\end{equation}
These are useful theoretical extensions for observables of an inverse.
Version 1.0 implements $g$ itself and the marker formula; it does not export
a separate observable-power API. Ordinary symbolic operations on a finite
\texttt{Normal} result must not be confused with automatically propagated
asymptotic error bounds.

\section{Existence, uniqueness, and convergence of the expansion}
\label{sec:convergence}

\subsection{A triangular formal equation}

Write $g(y)=z(1+U)$ and define
\begin{equation}
 E(U,z,L)=(1+U)^p-1+
 \sum_i z^{\delta_i}(1+U)^{p+\delta_i}
 B_i\bigl(L+\log(1+U)\bigr).
 \label{eq:normalized-equation}
\end{equation}
The inverse equation is $E=0$, and $U$ has strictly positive valuation.
At any new support value $\sigma$, the unknown coefficient of
$z^\sigma$ enters as $pC_\sigma$ plus an expression determined by lower
weights. Hence every coefficient is determined by one division by $p$.
This proves formal existence recursively and formal uniqueness. It also
provides a second coefficient algorithm independent of Lagrange inversion.

More explicitly, if $U$ and $V$ are two solutions with the same leading
constant, then
\[
 E(U)-E(V)=(U-V)(p+\text{positive-order terms}).
\]
The second factor is a unit, so $U=V$. The right inverse is also a left
inverse: composing it with $f$ gives the solution asymptotic to $x$ of
$f(X)=f(x)$, and the same nonzero-leading-factor argument forces $X=x$.

\subsection{An analytic proof using finitely many small parameters}

A formal series alone does not establish an asymptotic statement. Here a
particularly concrete analytic realization is available.

\begin{theorem}[Convergence for a finite power--log model]
For fixed real parameters satisfying \eqref{eq:model}, the sum
\eqref{eq:master}, including its logarithmic polynomials, is absolutely
convergent for all sufficiently small positive $y$ and equals the positive
inverse selected in Section~\ref{sec:branch}.
\end{theorem}
\begin{proof}
Write $w=1+U$ and $B_i(v)=\sum_{r=0}^{d_i}b_{ir}v^r$. Expansion in the
logarithmic argument gives
\[
 B_i(L+\log w)=\sum_{j=0}^{d_i}L^j A_{ij}(\log w),\qquad
 A_{ij}(v)=\sum_{r=j}^{d_i}b_{ir}\binom rj v^{r-j}.
\]
Introduce independent complex parameters $T_{ij}$ and the equation
\begin{equation}
 \mathcal E(w,T)=w^p-1+
 \sum_{i,j}T_{ij}w^{p+\delta_i}A_{ij}(\log w)=0.
 \label{eq:analytic-parameters}
\end{equation}
On a fixed complex disk about $w=1$ of radius less than one, select the
analytic logarithm with $\log1=0$. All powers and logarithms in this
formula are analytic there. At $(w,T)=(1,0)$,
$\mathcal E=0$ and $\partial_w\mathcal E=p\ne0$.
The analytic implicit function theorem gives an analytic function $w=W(T)$
on a polydisc about zero, with $W(0)=1$.

Now specialize
\begin{equation}
 T_{ij}=z^{\delta_i}(\log z)^j.
 \label{eq:physical-parameters}
\end{equation}
Every component tends to zero as $z\to0^+$, so the convergent Taylor series
of $W$ applies. Its value is real and close to one, hence positive, because
all coefficients and the physical parameters are real. It solves the
original inverse equation and therefore agrees with the unique local
positive inverse.

The Taylor series is absolutely convergent in a smaller polydisc. Group
its monomials by the total number $n_i$ of parameters with first index $i$.
For each such multi-index there are only finitely many distributions among
the logarithmic labels $j$. Absolute convergence permits regrouping.
Marker coefficient extraction, or formal uniqueness of
\eqref{eq:normalized-equation}, identifies the grouped coefficient with
\eqref{eq:coefficient}. This proves the claimed convergence and equality.
\end{proof}

The theorem does not give an ordinary radius of convergence in $y$ for a
Taylor series: logarithms and irrational powers are still present. It gives
a convergent analytic function of finitely many small composite parameters.
Nor does it give a uniform neighborhood as coefficients or gaps vary.
Fixed-parameter asymptotics is the assertion proved here.

\subsection{Newton iteration and order doubling}

A useful alternative implementation would iterate
\begin{equation}
 U_{k+1}=U_k-\frac{E(U_k)}{\partial_U E(U_k)},\qquad U_0=0,
 \label{eq:newton}
\end{equation}
with all operations truncated in the positive-weight algebra. Since
$\partial_U E=p+$ positive-order terms, division is by a unit.
If $U_*$ is the formal solution and
$\val(U_k-U_*)\ge s>0$, a Taylor expansion in the displacement shows
\[
 \val(U_{k+1}-U_*)\ge2s.
\]
Indeed, the linear displacement is canceled by the Newton quotient, while
all coefficients of the remaining quadratic and higher powers have
nonnegative valuation. Initially $s\ge\min_i\delta_i$.
This is a formal order-doubling statement, not a floating-point convergence
claim. Version 1.0 uses the explicit multi-index formula instead: it is easy
to audit, exposes the support, and provides a convenient boundary for the
remainder. Its residual checker implements the sparse composition needed
for an eventual Newton engine, but is not itself a Newton solver.

\section{Truncation and rigorous remainder semantics}
\label{sec:remainder}

\subsection{A finite-boundary theorem}

Suppose the requested cutoff in the target variable is $q>1/p$. Define
\begin{equation}
 H=pq-1,\qquad
 I_H=\{\boldsymbol n\in\N^m:\sigma(\boldsymbol n)<H\}.
 \label{eq:cutoff}
\end{equation}
This is a finite downward-closed set. Its one-step outer boundary is
\begin{equation}
 \partial I_H=\{\boldsymbol n+e_i:\boldsymbol n\in I_H,\ 1\le i\le m\}
                 \setminus I_H.
 \label{eq:boundary}
\end{equation}
The boundary need not be a minimal generating set. It is nevertheless
finite and every excluded multi-index dominates at least one boundary
index componentwise: follow a path from zero to it, increasing one
coordinate at a time, and take the first point leaving $I_H$.

If $m>0$, put
\begin{equation}
 \sigma_* =\min_{\nu\in\partial I_H}\sigma(\nu),\qquad
 k_* =\max_{\substack{\nu\in\partial I_H\\\sigma(\nu)=\sigma_*}}D(\nu),
 \qquad \beta_*=(1+\sigma_*)/p.
 \label{eq:boundary-data}
\end{equation}
Let $G_q$ be the sum in \eqref{eq:master} over $I_H$.

\begin{theorem}[Remainder from a finite boundary]
For the finite exact model,
\begin{equation}
 g(y)-G_q(y)=\OO\left(y^{\beta_*}(1+|\log y|)^{k_*}\right).
 \label{eq:boundary-remainder}
\end{equation}
In particular $\beta_*\ge q$. If $m=0$, the inverse is exactly $(y/a)^{1/p}$
and the remainder is zero.
\end{theorem}
\begin{proof}
Use the analytic parameter representation from the preceding section.
By a Cauchy estimate on a smaller polydisc, the absolute values of grouped
Taylor coefficients admit a convergent geometric majorant in the variables
\[
 u_i=Cz^{\delta_i}(1+|L|)^{d_i},
\]
for some fixed $C$. The finite number of logarithmic labels per block is
absorbed by decreasing the polydisc radius. The contribution of all terms
whose grouped multi-index is outside $I_H$ is therefore bounded by a
constant times
\[
 \sum_{\nu\in\partial I_H}u^\nu
       \prod_i(1-u_i)^{-1}
\]
for sufficiently small $z$. This can overcount terms, which is harmless for
an absolute bound. The geometric factors remain bounded because $u_i\to0$.
Undoing the definition of $u_i$ gives an error in $w=g/z$ bounded by a finite
sum of $z^{\sigma(\nu)}(1+|L|)^{D(\nu)}$.

Among a finite set of such bounds, the smallest power dominates; at that
power the largest logarithmic degree dominates. Equation
\eqref{eq:dominance} absorbs every larger power, regardless of its finite
logarithmic degree. Multiplication by $z$ and the fixed substitution
$z=(y/a)^{1/p}$ prove \eqref{eq:boundary-remainder}.
\end{proof}

\paragraph{What this bound does and does not say.}
$\beta_*$ is the first \emph{possible} omitted support value. It is not
necessarily the first nonzero omitted coefficient, because resonances or
special parameter values can cancel it. The reported bound remains valid
in that case, but may not be sharp. Similarly $k_*$ is a safe degree bound,
not a guarantee that its leading logarithmic coefficient is nonzero.
No numerical constant or explicit interval $0<y<y_0$ is computed by this
theorem or by the package.

The inert object \texttt{PowerLogRemainder[y,b,k]} means precisely
$\OO(y^b(1+|\log y|)^k)$ as $y\to0^+$. It is not an ordinary expression to
which power-series arithmetic should be applied. In particular, taking
\texttt{Normal} discards error metadata; subsequent symbolic operations on
that expression do not automatically propagate the error.

\subsection{Why input precision has to be transported}

Often a complicated elementary function is represented by a finite forward
jet rather than by an exact finite formula. Suppose
\begin{equation}
 F(x)=f(x)+R(x),\qquad
 R(x)=\OO(x^\rho M(x)^k),\quad
 R'(x)=\OO(x^{\rho-1}M(x)^k),
 \label{eq:input-remainder}
\end{equation}
where $M(x)=1+|\log x|$, $\rho>p$, and $k\ge0$ is an integer.

\begin{theorem}[Forward-to-inverse precision transport]
Under \eqref{eq:input-remainder}, $F$ has a positive local inverse and
\begin{equation}
 F^{-1}(y)-f^{-1}(y)=
 \OO\left(y^{(\rho-p+1)/p}(1+|\log y|)^k\right).
 \label{eq:precision-transport}
\end{equation}
\end{theorem}
\begin{proof}
The bounds imply $F(x)\sim ax^p$ and $F'(x)\sim apx^{p-1}$, so its
positive inverse exists and is asymptotic to $z=(y/a)^{1/p}$.
Let $u=F^{-1}(y)$ and $v=f^{-1}(y)$; then both are asymptotic to $z$.
The residual at $v$ is $F(v)-y=R(v)=\OO(z^\rho M(z)^k)$.
On the interval joining $u$ and $v$, the derivative of $F$ is asymptotic,
uniformly along this shrinking comparable interval, to $apz^{p-1}$.
The mean value theorem therefore gives
$u-v=\OO(z^{\rho-p+1}M(z)^k)$. Converting from $z$ to $y$ gives the result.
\end{proof}

The derivative condition is a hypothesis, not a consequence of the value
bound. For example, $R(x)=x^2\sin(x^{-2})$ is $\OO(x^2)$, but
$R'(x)=2x\sin(x^{-2})-2x^{-1}\cos(x^{-2})$ can destroy the monotonicity of
$x+R(x)$. An independently proved monotonicity hypothesis can sometimes
replace derivative control for value-error transport; the package uses the
simple differentiated-remainder contract to avoid silently assuming it.

The option \texttt{InputRemainder -> \{rho,k\}} records these hypotheses.
The implementation requires $\rho$ to exceed every supplied forward power,
and rejects a requested cutoff larger than
\begin{equation}
 q_{\rm input}=(\rho-p+1)/p.
 \label{eq:input-cap}
\end{equation}
This is a deliberately conservative interface. Without more detailed
logarithmic precision at the boundary, coefficients at or beyond that
power are not justified. The final remainder is the larger of the finite
model's truncation error and the transported input error: compare powers
first, and take the larger logarithmic degree when powers agree.

\paragraph{A correct elementary-function workflow.}
For $F(x)=e^x-1+x^{\sqrt2}$, take the Taylor polynomial of the analytic part
through $x^4$ and retain the exact irrational-power term:
\begin{lstlisting}
jet = Normal[Series[Exp[x] - 1, {x, 0, 4}]] + x^Sqrt[2];
r = RealInverseAsymptotic[jet, {x, 0}, {y, 5},
      InputRemainder -> {5, 0}];
\end{lstlisting}
Taylor's theorem supplies both bounds in \eqref{eq:input-remainder}.
Dropping an \texttt{O} term with \texttt{Normal} and omitting
\texttt{InputRemainder} would instead ask the package to invert the finite
polynomial model \emph{exactly}. That is a different mathematical input.

\section{The logarithmic example to all orders}
\label{sec:logexample}

Here $a=p=\delta_1=1$ and $B_1(L)=1+L$. Formula
\eqref{eq:master} takes the simple form
\begin{equation}
 g_1(y)=y+\sum_{n\ge1}y^{n+1}P_n(\log y),\qquad
 P_n(L)=\frac{(-1)^n}{n!}
          \prod_{j=n+2}^{2n}(\D+j)(1+L)^n.
 \label{eq:log-all-orders}
\end{equation}
This also follows immediately from \eqref{eq:near-id} by differentiating
$y^{2n}(1+\log y)^n$. The two forms are useful cross-checks of the operator
normalization.

\begin{proposition}[Highest logarithms are Catalan coefficients]
$P_n$ has degree exactly $n$, and
\begin{equation}
 [L^n]P_n(L)=(-1)^n\frac{1}{n+1}\binom{2n}{n}.
 \label{eq:catalan-leading}
\end{equation}
Consequently, after $n=N$ in \eqref{eq:log-all-orders}, the remainder is
\begin{equation}
 \OO\bigl(y^{N+2}(1+|\log y|)^{N+1}\bigr),
 \label{eq:log-tail}
\end{equation}
and the displayed power--log order is attained by the next block.
\end{proposition}
\begin{proof}
A derivative lowers polynomial degree. To obtain the coefficient of $L^n$,
select the scalar part in every operator factor. This gives
\[
 \frac{(-1)^n}{n!}\prod_{j=n+2}^{2n}j
 =(-1)^n\frac{(2n)!}{n!(n+1)!}\ne0.
\]
The boundary index is $N+1$, of weight $N+1$ and logarithmic degree $N+1$,
so the remainder theorem gives \eqref{eq:log-tail}. Its first omitted
coefficient has the nonzero leading logarithm just calculated.
\end{proof}

The first four $P_n$ appear in \eqref{eq:first-answer}. Two more are
\begin{align}
 P_5(L)={}&-42L^5-\frac{3727}{12}L^4-\frac{5365}{6}L^3
           -1256L^2-\frac{5177}{6}L-\frac{2789}{12},\label{eq:p5}\\
 P_6(L)={}&132L^6+\frac{5981}{5}L^5+\frac{52937}{12}L^4
          +\frac{16979}{2}L^3\nonumber\\
         &+9007L^2+\frac{30019}{6}L+\frac{22769}{20}.
 \label{eq:p6}
\end{align}
The coefficient table in the archive is generated independently by SymPy.
Its highest logarithms agree with \eqref{eq:catalan-leading}, and direct
sparse substitution cancels the normalized residual at every retained
power.

\subsection{A residual that identifies the correct error scale}

Let $G_2(y)=y-y^2(1+L)$. Substitution and expansion give
\begin{equation}
 \frac{f_1(G_2(y))}{y}-1
 =-y^2(2L^2+5L+3)+\OO(y^3(1+|L|)^3).
 \label{eq:first-residual}
\end{equation}
Since $f_1'\to1$, the forward residual and the inverse error have the same
leading magnitude and opposite signs under this normalization. This gives
an independent explanation of the missing third block and the invalidity
of a bare $\OO(y^3)$ in \eqref{eq:not-o-y3}.

\subsection{An explicit a posteriori error bound}

For any positive approximation $G$ and any $y>0$, the global derivative
bound \eqref{eq:global-slope} implies
\begin{equation}
 \abs{G-g_1(y)}\le\frac{\abs{f_1(G)-y}}{1-2e^{-5/2}}.
 \label{eq:global-certificate}
\end{equation}
This follows directly from the mean value theorem, not from a formal
series identity. When the residual is evaluated with rigorous interval
arithmetic, it can be made into a rounding-certified numerical enclosure.
Ordinary high-precision evaluation is evidence, not a substitute for such
an interval computation. The supplied numerical tables use high precision
and are explicitly diagnostics rather than interval certificates.

\section{The irrational-power example and its scalar reduction}
\label{sec:irrational}

The entire family
\begin{equation}
 f_\alpha(x)=x+x^\alpha,\qquad\alpha>1,
 \label{eq:alpha}
\end{equation}
is globally increasing on $(0,\infty)$ because
$f_\alpha'(x)=1+\alpha x^{\alpha-1}>1$. Its endpoint limits are zero and
infinity. Thus the positive inverse is unambiguous.

\begin{theorem}[Generalized Fuss--Catalan inverse coefficients]
For fixed $\alpha>1$,
\begin{equation}
 g_\alpha(y)=y+\sum_{n\ge1}
       \frac{(-1)^n}{n}\binom{\alpha n}{n-1}
       y^{1+n(\alpha-1)}.
 \label{eq:alpha-series}
\end{equation}
The binomial coefficient here is the finite product
\begin{equation}
 \frac1n\binom{\alpha n}{n-1}
 =\frac1{n!}\prod_{j=0}^{n-2}(\alpha n-j)
 =\frac1{1+n(\alpha-1)}\binom{\alpha n}{n}.
 \label{eq:binomial-product}
\end{equation}
No rationality assumption on $\alpha$ is needed.
\end{theorem}
\begin{proof}
Apply \eqref{eq:near-id} to $h(y)=y^\alpha$ and differentiate
$y^{\alpha n}$ exactly $n-1$ times. This gives the finite product in
\eqref{eq:binomial-product} and the power in \eqref{eq:alpha-series}.
The other two forms are elementary rearrangements of the same product.
\end{proof}

For $\alpha=\sqrt2$, the first coefficients are $-1$, $\sqrt2$, and
$-(6-\sqrt2)/2$, reproducing the requested result.
After correction $n=N$, the remainder is
\begin{equation}
 \OO\left(y^{1+(N+1)(\alpha-1)}\right).
 \label{eq:alpha-remainder}
\end{equation}
The package call that retains exactly the leading term and three corrections
in this example is
\begin{lstlisting}
r = RealInverseAsymptotic[
  x + x^Sqrt[2], {x, 0}, {y, 4 Sqrt[2] - 3}];
Normal[r]
\end{lstlisting}
A power cutoff, rather than a count of ordinary integer powers, is the
relevant request.

\subsection{A one-variable analytic generating function}

Put $t=y^{\alpha-1}$ and $g_\alpha(y)=yV(t)$. The defining equation becomes
\begin{equation}
 V+tV^\alpha=1,\qquad V(0)=1.
 \label{eq:V}
\end{equation}
It is analytic in $(V,t)$ near $(1,0)$, so $V$ has an ordinary convergent
Taylor series in $t$. Equation \eqref{eq:alpha-series} is exactly that
Taylor series multiplied by $y$. This simple reduction is excellent for a
single gap; with several incommensurable gaps there need not be a single
parameter producing an integer power grid.

In fact the radius of convergence of $V$ is
\begin{equation}
 R_\alpha=\frac{(\alpha-1)^{\alpha-1}}{\alpha^\alpha}.
 \label{eq:radius}
\end{equation}
To verify it directly, let $c_n$ be the positive number in
\eqref{eq:binomial-product}. Its gamma-function form and Stirling's formula
give
\begin{equation}
 c_n\sim\frac{\sqrt\alpha}{\sqrt{2\pi}(\alpha-1)^{3/2}}
 n^{-3/2}R_\alpha^{-n}.
 \label{eq:cn-asymptotic}
\end{equation}
The root test proves \eqref{eq:radius}, and the $n^{-3/2}$ factor also proves
absolute convergence at $|t|=R_\alpha$. As a consistency check, the first
real critical point of \eqref{eq:V} is
$V=\alpha/(\alpha-1)$, $t=-R_\alpha$, obtained by setting its derivative
with respect to $V$ to zero.

This result is not uniform as $\alpha\downarrow1$: the gap
$\alpha-1$ tends to zero, the asymptotic hierarchy collapses, and at
$\alpha=1$ the inverse is $y/2$, not $y+o(y)$. The software consequently
requires fixed numerical algebraic exponents instead of pretending to
order symbolic gaps uniformly in a parameter.

\section{Further examples and reductions}
\label{sec:examples}

\subsection{Nonunit leading powers}

Consider
\[
 f(x)=3x^2\{1+x(1+\log x)\}.
\]
With $z=\sqrt{y/3}$ and $L=\log z$,
\begin{equation}
 g(y)=z-\frac12z^2(1+L)
       +\frac18z^3(5L^2+12L+7)
       +\OO\bigl(y^2(1+|\log y|)^3\bigr).
 \label{eq:leading-two}
\end{equation}
The first two corrections follow from \eqref{eq:first-two-general} with
$p=2$, $\delta=1$. Notice both the factor $p^{-N}$ in the coefficient and
the change $L=\log(y/3)/2$. Forgetting either produces an incorrect inverse.

\subsection{A resonance and an exact cancellation}

For $f(x)=x+x^2+x^3$, take gaps $1$ and $2$. At inverse power $y^3$, the
index $(2,0)$ contributes $2$ and $(0,1)$ contributes $-1$, leaving $1$.
At power $y^4$, the index $(3,0)$ contributes $-5$ and $(1,1)$ contributes
$5$, so the entire block disappears. A power cutoff must be applied to
weights, and terms of equal weight must be combined before output.
A remainder bound based on a canceled boundary block is still valid, but
is not sharp.

\subsection{Several incommensurable gaps and logarithms}

The model
\begin{equation}
 f(x)=x+x^{\sqrt2}(1+\log x)+2x^{\sqrt3}
 \label{eq:mixed}
\end{equation}
has gaps $\sqrt2-1$ and $\sqrt3-1$.
The weights include each gap, twice the first gap, and their sum; logarithmic
polynomials occur at some but not all weights. Formula
\eqref{eq:cross-term} supplies the mixed correction. The package generates
all weights below the requested cutoff without rationalizing either gap.
An independent exact residual test for this example is included.

\subsection{Other finite endpoints and inversion at infinity}

For a right-sided branch at $x=x_0$, $y=y_0$, introduce $u=x-x_0>0$ and
$v=y-y_0$. If $F(x_0+u)-y_0$ has the required form, invert it and add $x_0$.
For a left-sided branch, choose $u=x_0-x>0$ and orient the target coordinate
so that its leading coefficient is positive. These changes of variables
are mathematical branch data; version 1.0 accepts the normalized endpoint
zero and does not choose them automatically.

If a positive function grows to infinity and
$\widetilde F(t)=1/F(1/t)$ has an admissible expansion at $t=0^+$, then
\[
 F^{-1}(y)=\frac1{\widetilde F^{-1}(1/y)}.
\]
The forward reciprocal must first be expanded with its correct remainder,
and inverse error must be transported through the final reciprocal.
A normal expression alone does not retain that precision information.

\section{Implementation and its exact contract}
\label{sec:implementation}

\subsection{Accepted and rejected inputs}

The main engine accepts a finite sum
\[
 f(x)=\sum_j x^{\alpha_j}P_j(\log x)
\]
whose smallest nonzero block is $a x^p$, $a,p>0$.
All exponents and the target cutoff must be exact real algebraic constants.
Logarithmic coefficients must be polynomials with nonnegative integer
logarithmic degree; their coefficients may contain fixed real parameters
when \texttt{Assumptions} proves reality and the leading coefficient's
positivity. Inexact \texttt{Real} input is rejected rather than silently
converted to exact rational data.

The parser intentionally rejects unexpanded analytic factors, symbolic
exponents, negative powers of $\log x$, iterated logarithms, exponential
sectors, and a leading logarithmic block. This is a supported-class boundary,
not a statement that such functions have no inverse expansion. See
Section~\ref{sec:extensions} for the required extensions.

The source and target must be distinct unassigned symbols, and the forward
expression must not contain the target symbol. User assumptions concern
fixed parameters only, not either asymptotic variable. Positivity of source
and target is already part of the selected branch. Use exact
\texttt{Sqrt[2]}, not a decimal approximation to it.

\subsection{Public functions and data}

\begingroup\small
\begin{longtable}{>{\raggedright\arraybackslash}p{0.38\textwidth}>{\raggedright\arraybackslash}p{0.56\textwidth}}
\toprule
Interface & Meaning\\\midrule\endhead
\texttt{RealInverseAsymptotic[}\newline\texttt{f,\{x,0\},\{y,q\}]} &
Construct the complete-block inverse expansion below the exclusive power cutoff $q$.\\
\texttt{Normal[result]} &
The finite expression only; discard its remainder metadata.\\
\texttt{result["Remainder"]} &
Zero for an exact monomial inverse, otherwise an inert \texttt{PowerLogRemainder} descriptor.\\
\texttt{result["Terms"]} &
Pairs $\{\beta,C\}$ meaning $(y/a)^\beta C$, with logarithms already substituted as $\log(y/a)/p$.\\
\texttt{result["Branch"]} &
The positive branch with $x/(y/a)^{1/p}\to1$.\\
\texttt{InverseResidual[result]} &
Independent normalized sparse-composition residual below the result's relative cutoff.\\
\texttt{InverseResidual[result,h]} &
The same calculation at a user-specified positive relative cutoff in the $z$ coordinate. It can expose the first nonzero omitted residual.\\
\texttt{PerturbativeInverse[}\newline\texttt{h,\{x\},\{y,n\},eps]} &
Degree-$n$ polynomial in a marker for $u+\texttt{eps}\,h(u)=y$, with $u=y$ at marker zero.\\
\bottomrule
\end{longtable}
\endgroup

The option \texttt{Assumptions} defaults to \texttt{True},
\texttt{InputRemainder} to \texttt{None}, and
\texttt{MaxMultiIndices} to \texttt{20000}. The last is a conservative resource
budget for traversal, boundary candidates, and sparse products; it is not
an exact count of final output blocks and not a global time limit.
A budget failure returns \texttt{Failure}, never a result pretending to have
reached the requested order.

Additional fields include \texttt{"LeadingPower"},
\texttt{"LeadingCoefficient"}, \texttt{"InputInterpretation"},
\texttt{"RelativeCutoff"}, \texttt{"RemainderPower"},
\texttt{"RemainderLogDegree"}, \texttt{"MultiIndexCount"}, and
\texttt{"BoundaryCount"}. Internal scaled terms are retained for exact
residual verification. The mathematical guarantee is conditional on the
fixed-parameter hypotheses; it is neither an interval certificate nor a
uniform parameter estimate.

\subsection{The algorithm in six stages}

First expand the finite input into monomial factors, replace
$\log x$ temporarily by a fresh symbol, and gather equal powers.
Prove reality of coefficient data and positivity of the constant leading
block. Second, normalize by $a x^p$ and compute the positive gaps.
Third, enumerate $I_H$ recursively with exact comparisons; prune a branch
as soon as its weight reaches $H$. Fourth, evaluate
\eqref{eq:coefficient} for every retained multi-index. Fifth, canonicalize
and gather equal weights, cancel zero polynomials, and convert to the target
variable. Sixth, compute the finite boundary and combine its error bound
with any transported input uncertainty.

The core polynomial loop is small:
\begin{lstlisting}
q = Expand[Times @@ MapThread[Power, {blocks, n}]];
Do[
  q = Expand[D[q, ell] + (1 + weight + p j) q],
  {j, 1, Total[n] - 1}
];
q = (-1)^Total[n] q /
    (p^Total[n] (Times @@ (Factorial /@ n)));
\end{lstlisting}
The surrounding validation and support arithmetic matter just as much as
this formula. \texttt{RootReduce} is used to canonicalize exact algebraic
weights, and comparisons must simplify to a definite truth value; there is
no approximate tie-breaking rule \cite{wlrootreduce}. The display layer
uses \texttt{ToRadicals} when available, without changing the exact keys.

If $\delta_{\min}=\min_i\delta_i$, the maximum total degree retained is
strictly smaller than $H/\delta_{\min}$. A coarse upper bound on the number
of retained indices is
\[
 \prod_i\left(1+\left\lfloor H/\delta_i\right\rfloor\right).
\]
The actual simplex enumeration is smaller, but high dimension, tiny gaps,
large logarithmic degrees, or complicated algebraic numbers can still make
a request expensive. These are mathematical and representation costs, not
evidence that a rational grid should be substituted for an irrational one.

\subsection{Independent residual verification}

For a computed finite $U$, the checker forms \eqref{eq:normalized-equation}
without using the Lagrange coefficient formula. Sparse convolution is used
for multiplication, and
\[
 (1+U)^r=\sum_{j\ge0}\binom rjU^j,\qquad
 \log(1+U)=\sum_{j\ge1}\frac{(-1)^{j+1}}jU^j
\]
are truncated as soon as their valuations exceed the cutoff. Polynomial
substitution $B_i(L+\log(1+U))$ is a finite Taylor expansion in its second
summand. This checks normalization, logarithmic differentiation, mixed
terms, and resonant cancellation by a different route.

The returned residual is
\[
 f(G_q(y))/y-1.
\]
If its relative cutoff is $H$, the corresponding cutoff for the unnormalized
forward residual is $1+H/p$. A zero residual below the cutoff confirms the
finite formal identity for the supplied finite model. It does not prove the
user's unspecified remainder bounds and does not assess floating-point
rounding. If the result was constructed from a forward jet, the check is
explicitly limited to that finite jet.

\subsection{Why not manufacture a rational \texttt{SeriesData} object?}

\texttt{SeriesData} stores powers on its documented index/denominator grid;
logarithmic factors can appear in its coefficients
\cite{wlseriesdata}. That is useful for ordinary series and many logarithmic
expansions, but it is not a faithful reason to round
$1+n(\sqrt2-1)$ to rational exponents. The package uses its own transparent
finite data object and explicit remainder descriptor instead. It also does
not overload big-O arithmetic that it has not implemented.

\section{Verification performed and reproducibility}
\label{sec:verification}

\subsection{What was actually executed}

The connected Wolfram evaluator was attempted, but its MCP endpoint returned
HTTP 404 before evaluating code. No local Wolfram kernel was available.
Accordingly, the native \texttt{.wlt} suite is supplied but is \emph{not}
reported as executed, and there is no claim of a tested Mathematica version.

The independent Python verification script was executed. It uses SymPy for
exact algebra and mpmath at 150 decimal digits for numerical comparisons.
It reports 18 successful checks, including exact sparse-composition tests
for the original logarithmic and irrational examples, mixed irrational gaps,
resonances, two logarithmic blocks, a nonunit leading power, a fractional
leading power, and a pure monomial. It independently compares the first
six logarithmic corrections with the derivative form of Lagrange inversion,
checks their signed Catalan leading coefficients, checks an ordinary
quadratic inverse against its square-root formula, and verifies the first
nonzero omitted logarithmic residual. The 24 numerical rows and 12 numerical
checks of \eqref{eq:global-certificate} are diagnostics, not certified
interval calculations.

A lexical check of balanced delimiters, strings, and nested comments was
also performed on the Wolfram sources. This detects a limited class of
source defects; it does not execute Wolfram semantics or certify the parser.
The limitations above are important even though independent mathematical
checks are substantial. They are recorded in
\texttt{verification/STATUS.md} and the machine-readable result file.

\subsection{Numerical behavior on the positive branch}

In Table~\ref{tab:lognumeric}, $m$ is the largest retained ordinary power
of $y$ for the logarithmic example, and the scaled error is
\[
 \frac{|G_m(y)-g_1(y)|}{y^{m+1}(1+|\log y|)^m}.
\]
The reference inverse was obtained by high-precision bisection on the
positive interval $[y/2,2y]$, which brackets the root for the displayed
arguments. Table~\ref{tab:irrnumeric} similarly uses $N$ correction terms
for the irrational example and divides the error by
$y^{1+(N+1)(\sqrt2-1)}$. The tables illustrate convergence and the proven
orders; they are not the proofs of those orders.

\begin{table}[htbp]
\centering\small
\begin{tabular}{rrrr}
\toprule
$y$ & $m$ & Absolute error & Error / stated scale \\
\midrule
$10^{-2}$ & 2 & $2.4202\times10^{-5}$ & 0.770327 \\
$10^{-2}$ & 3 & $1.8127\times10^{-6}$ & 1.02936 \\
$10^{-2}$ & 5 & $1.2899\times10^{-8}$ & 2.33140 \\
$10^{-2}$ & 7 & $1.0791\times10^{-10}$ & 6.20810 \\
$10^{-4}$ & 2 & $1.2685\times10^{-10}$ & 1.21677 \\
$10^{-4}$ & 3 & $2.4098\times10^{-13}$ & 2.26395 \\
$10^{-4}$ & 5 & $1.1512\times10^{-18}$ & 10.3744 \\
$10^{-4}$ & 7 & $6.6354\times10^{-24}$ & 57.3572 \\
$10^{-8}$ & 2 & $5.8954\times10^{-22}$ & 1.56309 \\
$10^{-8}$ & 3 & $2.4782\times10^{-28}$ & 3.38338 \\
$10^{-8}$ & 5 & $5.8498\times10^{-41}$ & 21.1747 \\
$10^{-8}$ & 7 & $1.6739\times10^{-53}$ & 160.652 \\
\bottomrule
\end{tabular}
\caption{Positive-branch numerical diagnostics for the logarithmic example. Computations use 150 decimal digits; displayed values are rounded.}
\label{tab:lognumeric}
\end{table}

\begin{table}[htbp]
\centering\small
\begin{tabular}{rrrr}
\toprule
$y$ & $N$ & Absolute error & Error / stated scale \\
\midrule
$10^{-2}$ & 1 & $2.5197\times10^{-4}$ & 1.14343 \\
$10^{-2}$ & 2 & $5.9671\times10^{-5}$ & 1.82407 \\
$10^{-2}$ & 3 & $1.5336\times10^{-5}$ & 3.15815 \\
$10^{-2}$ & 4 & $4.1555\times10^{-6}$ & 5.76451 \\
$10^{-4}$ & 1 & $6.6314\times10^{-8}$ & 1.36556 \\
$10^{-4}$ & 2 & $2.3627\times10^{-9}$ & 2.20788 \\
$10^{-4}$ & 3 & $9.0972\times10^{-11}$ & 3.85768 \\
$10^{-4}$ & 4 & $3.6835\times10^{-12}$ & 7.08808 \\
$10^{-8}$ & 1 & $3.3324\times10^{-15}$ & 1.41310 \\
$10^{-8}$ & 2 & $2.6235\times10^{-18}$ & 2.29095 \\
$10^{-8}$ & 3 & $2.2302\times10^{-21}$ & 4.01029 \\
$10^{-8}$ & 4 & $1.9927\times10^{-24}$ & 7.37862 \\
\bottomrule
\end{tabular}
\caption{Positive-branch numerical diagnostics for the irrational example. Computations use 150 decimal digits; displayed values are rounded.}
\label{tab:irrnumeric}
\end{table}


\subsection{Running the supplied checks}

In a Wolfram environment, run
\begin{lstlisting}
wolframscript -file Tests/run-tests.wls
\end{lstlisting}
or evaluate the equivalent \texttt{TestReport} call from a notebook.
The supplied MUnit suite contains 36 tests covering expected coefficients,
residuals, remainder metadata, symbolic real coefficients, precision caps,
unsupported inputs, and resource-limit behavior. Results from this suite
must be obtained in the user's own kernel; no pass count is asserted here.

For the independent checks,
\begin{lstlisting}
python verification/verify.py --output verification
\end{lstlisting}
regenerates the JSON report and coefficient table. The script explicitly
identifies itself as a reference implementation, not a Wolfram interpreter.
The examples file contains the two original calls, a mixed model, a nonunit
leading model, a forward-jet example, and a high-precision numerical
comparison. The archive contains all LaTeX inputs needed to rebuild this
PDF; it does not depend on repository macros or on redistributing fonts.

\section{Boundaries of the theory implemented here}
\label{sec:extensions}

\subsection{A small logarithmic correction with zero power gap}

For $f(x)=x+x/\log x$, the correction is relatively small, but its power gap
is zero. All corrections can occur at the same power $y$, with increasingly
negative powers of $\log y$. The finite region $I_H$ would no longer be
finite in the power grading used here. The appropriate coordinate is
$t=1/\log y$, not another positive power of $y$. With $g=yW$,
\[
 W\left(1+\frac{t}{1+t\log W}\right)=1,
\]
which is analytic near $(W,t)=(1,0)$ and can be reversed as a Taylor series
in $t$. The main engine correctly rejects this input rather than attempting
to enumerate infinitely many zero-weight contributions.

\subsection{A leading logarithm, iterated logarithms, and flat sectors}

A leading behavior such as $x^p(-\log x)^b$ requires a different dominant
coordinate, often with a Lambert-$W$ core or an implicit logarithmic
normalization. The general identity \eqref{eq:general-core} remains useful,
but a theorem controlling its actual asymptotic grading must accompany it.
Leading logarithms are not automatically reduced by version 1.0.

Iterated logarithms or rational functions of $\log x$ require a larger
coefficient algebra and a second truncation policy. Exponentially small
terms such as $e^{-1/x}$ require an exponential sector grading. They are
not ordinary power--log coefficients, and silently dropping them would lose
information beyond all algebraic orders. The marker utility can generate
finite differential expressions for many such perturbations, but it does
not assert that setting the marker equal to one has an established
asymptotic remainder. This is the precise distinction between a formula
generator and a transseries engine with a proved support contract.

\subsection{Symbolic exponents and nonuniform limits}

When an exponent depends on a parameter, the ordering of support values can
change at resonances, and a positive gap may approach zero. A symbolic
exponent engine would need to partition parameter space into cells with
proved ordering relations, handle boundary cells separately, and carry
uniformity conditions into error bounds. The present exact algebraic
restriction avoids making unsupported decisions about such orderings.
The mathematics of the master formula still applies to any fixed positive
real gaps, but that does not mean a symbolic implementation can compare all
parameter-dependent weights without additional information.

\subsection{What a larger implementation should add next}

The most direct extension is a sparse Newton solver for the same supported
class, using the already implemented composition operations and a unit
reciprocal routine. It should be benchmarked against the explicit formula
on high-order and many-gap examples. A second extension is structured
forward-jet arithmetic, so elementary functions can be expanded together
with their differentiated remainder contracts rather than relying on a
manual \texttt{InputRemainder} declaration. Larger logarithmic and
exponential coefficient algebras should come with their own exact support
comparators and error semantics, not merely additional parser patterns.
Finally, interval evaluation of the residual and a proved derivative bound
would turn the a posteriori theory into computable numerical certificates.
None of these additions is falsely advertised as present in version 1.0.

\section{Relationship to the supplied materials and main conclusions}
\label{sec:relation}

The supplied repository directory contains a canonical transseries volume
and a combinatorial companion. The group README was read, and the relevant
phase-coordinate, multiplicative-sector, and scaled-reversion portions of
the companion source were inspected \cite{reporeadme,repocompanion}.
In particular, the companion's labels
\texttt{t1:thm:phase}, \texttt{t1:thm:multi}, and
\texttt{t2:prop:scaled} explain the operator and marker methodology.
The README explicitly warns against conflating formal identities, analytic
convergence, and the scope of formalized results.

The multi-megabyte canonical source could not be read through the available
repository-file endpoint; this report does not claim a complete review of
that volume. It instead gives self-contained proofs of every specialized
inversion and remainder result it uses. No new Lean formalization is
claimed, and the kernel checks of this Wolfram package are separate from
any formalization claims made by the repository.

The classical Lagrange identity is not new. The contribution here is its
explicit positive-endpoint specialization with arbitrary positive gaps,
nonunit monomial cores, a directly enumerable exact support, complete
logarithmic blocks, a finite-boundary remainder law, guarded forward
precision, and an auditable package interface tailored to the question.
The analysis separates three assertions that must never be confused:
coefficient correctness, realization as an asymptotic expansion of a
specified real branch, and successful execution in a particular CAS.
The first two are proved for the stated class and checked independently;
the third still requires the supplied suite to be run in an actual Wolfram
kernel.

The practical answer is consequently not a request to persuade
\texttt{InverseFunction} to guess the intended branch. Select the positive
leading root, compute finite power--log blocks in their natural support,
retain the logarithms in the remainder, and verify the forward residual.
Both original examples, and finite mixtures substantially more complicated
than either one, fit this construction.

\appendix
\section{A compact marker recipe}

For the specific near-identity class $f(x)=x+h(x)$, the following minimal
recipe is useful even independently of the full package:
\begin{lstlisting}
ClearAll[nearIdentityPolynomial];
nearIdentityPolynomial[h_, x_Symbol, y_Symbol, n_Integer] /; n >= 0 :=
  y + Total[Table[
    (-1)^j D[(h /. x -> y)^j, {y, j - 1}]/j!,
    {j, 1, n}
  ]];

nearIdentityPolynomial[x^2 (1 + Log[x]), x, y, 4] // Expand
\end{lstlisting}
The table uses numerical integer loop indices before taking the derivatives.
For the original logarithmic perturbation, the result has the remainder in
\eqref{eq:log-tail}; for a general $h$, this code alone does not establish
an asymptotic grading. The package's \texttt{PerturbativeInverse} keeps an
explicit marker to make that qualification visible.

\section{An exact residual bracket lemma}

\begin{lemma}[Existence and an error enclosure from a residual]
Let $F$ be continuously differentiable on $[A,B]$, with $F'\ge m>0$ there.
Let $x_0\in(A,B)$ and $r=F(x_0)-y$. If
\[
 [x_0-|r|/m,\ x_0+|r|/m]\subseteq[A,B],
\]
then $F(x)=y$ has a unique solution in that interval, and
$|x-x_0|\le |r|/m$.
\end{lemma}
\begin{proof}
Integrating the derivative lower bound gives
$F(x_0+|r|/m)\ge F(x_0)+|r|\ge y$ and
$F(x_0-|r|/m)\le F(x_0)-|r|\le y$.
The intermediate value theorem supplies a root; strict monotonicity gives
uniqueness. The stated interval is the error enclosure.
\end{proof}

This lemma establishes existence as well as an error bound, unlike a
mean-value estimate that presupposes a nearby root. For the first example
global existence has already been proved, so
\eqref{eq:global-certificate} is available directly. For general models,
an explicit interval and derivative bound would be extra computable inputs
to a numerical certification layer.

\section{Source and verification manifest}

\begingroup\small
\begin{longtable}{>{\raggedright\arraybackslash}p{0.47\textwidth}>{\raggedright\arraybackslash}p{0.47\textwidth}}
\toprule
Archive member & Purpose\\\midrule\endhead
\path{Kernel/RealInverseAsymptotics.wl} & Main package and independent formal residual arithmetic.\\
\path{Kernel/init.m} & Conventional package loader.\\
\path{Examples/examples.wl} & Reproducible Wolfram calls and numerical diagnostic.\\
\path{Tests/RealInverseAsymptotics.wlt} & 36 native MUnit regression tests; not executed in this session.\\
\path{Tests/run-tests.wls} & Native test runner with an exit status.\\
\path{verification/verify.py} & Independently executed exact and numerical reference checks.\\
\path{verification/verification_results.json} & Recorded Python results, versions, coefficients, and numerical values.\\
\path{verification/STATUS.md} & Explicit boundary between executed checks and unexecuted native tests.\\
\path{article/real_inverse_asymptotics.tex} & Main manuscript source.\\
\path{article/numeric_log.tex}, \path{numeric_irr.tex} & Generated table inputs used by the manuscript.\\
\path{article/real_inverse_asymptotics.pdf} & Compiled, visually inspected report.\\
\path{README.md}, \path{LICENSE} & Usage, build instructions, scope, and licensing.\\
\bottomrule
\end{longtable}
\endgroup

\begingroup\small
\begin{thebibliography}{20}
\bibitem{question}
V. Reshetnikov, \emph{How to get an asymptotic of the real-valued branch of
the inverse function?} Mathematica Stack Exchange, question 236367,
11 December 2020, updated 13 December 2020. User-supplied archive and public
page consulted 7 September 2026.
\url{https://mathematica.stackexchange.com/questions/236367/}.

\bibitem{gessel}
I. M. Gessel, \emph{Lagrange inversion}, Journal of Combinatorial Theory,
Series A \textbf{144} (2016), 212--249.
\url{https://arxiv.org/abs/1609.05988}.

\bibitem{edgarbegin}
G. A. Edgar, \emph{Transseries for beginners}, arXiv:0801.4877,
2008, revised 2009. Expository background on ordered transseries,
differentiation, and composition.
\url{https://arxiv.org/abs/0801.4877}.

\bibitem{edgarcomp}
G. A. Edgar, \emph{Transseries: Composition, Recursion, and Convergence},
arXiv:0909.1259, 2009. Background on compositional inversion and its formal
setting; the specialized proofs in this report are independent.
\url{https://arxiv.org/abs/0909.1259}.

\bibitem{reporeadme}
V. Reshetnikov / ProveIt project, \emph{Series and transseries}, group README,
consulted 7 September 2026; retrieved blob
\texttt{516aaa30e84cb0036a00f803cb6e8e3026febb8e}.
\url{https://github.com/VladimirReshetnikov/ProveIt/tree/main/Analysis/FabiusFunction/docs/semi-formalized-research-frontiers/drafts/series-and-transseries}.

\bibitem{repocompanion}
V. Reshetnikov / ProveIt formal-development synthesis,
\emph{Combinatorial Transseries and Their Inverses}, consolidated companion
volume, 4 September 2026. Relevant source sections on phase-coordinate
inversion and scaled local reversion; retrieved source blob
\texttt{05c56ce5ee60788a67f365110b62810550cba817}.
\url{https://github.com/VladimirReshetnikov/ProveIt/blob/main/Analysis/FabiusFunction/docs/semi-formalized-research-frontiers/drafts/series-and-transseries/Combinatorial_Transseries_Inverses/Combinatorial_Transseries_Inverses.tex}.

\bibitem{wlseriesdata}
Wolfram Research, \emph{SeriesData}, Wolfram Language documentation,
consulted 7 September 2026.
\url{https://reference.wolfram.com/language/ref/SeriesData.html}.

\bibitem{wlseries}
Wolfram Research, \emph{Series}, Wolfram Language documentation,
consulted 7 September 2026.
\url{https://reference.wolfram.com/language/ref/Series.html}.

\bibitem{wlinverseseries}
Wolfram Research, \emph{InverseSeries}, Wolfram Language documentation,
consulted 7 September 2026.
\url{https://reference.wolfram.com/language/ref/InverseSeries.html}.

\bibitem{wlinversefunction}
Wolfram Research, \emph{InverseFunction}, Wolfram Language documentation,
consulted 7 September 2026.
\url{https://reference.wolfram.com/language/ref/InverseFunction.html}.

\bibitem{wlpower}
Wolfram Research, \emph{Power}, Wolfram Language documentation,
consulted 7 September 2026.
\url{https://reference.wolfram.com/language/ref/Power.html}.

\bibitem{wlrootreduce}
Wolfram Research, \emph{RootReduce} and \emph{Algebraics}, Wolfram Language
documentation, consulted 7 September 2026.
\url{https://reference.wolfram.com/language/ref/RootReduce.html};
\url{https://reference.wolfram.com/language/ref/Algebraics.html}.
\end{thebibliography}
\endgroup
\end{document}
